\documentclass{article}
\usepackage[UTF8]{ctex}
\usepackage{amsmath,amssymb,amsfonts}
\usepackage{enumitem}
\usepackage{geometry}
\geometry{a4paper, margin=1in}
% === 错题本排版优化（不参与题目标识） ===
\setlength{\parindent}{0pt}
\setlength{\parskip}{0.35em plus 0.10em minus 0.08em}
\setlength{\emergencystretch}{3em}
\setlist{itemsep=0.25em, topsep=0.40em, parsep=0pt}
\allowdisplaybreaks[2]
% === 错题本排版优化结束 ===
\begin{document}
\section*{郑州大学2025年高等代数真题}
\section*{119．郑州大学2025年高等代数试卷}
一．填空题．每题5分，共30分．\\

\subsection{题目1}
1．\(A\) 是 3 阶矩阵，\(A^{*}\) 为 \(A\) 的伴随矩阵，若 \(|A|=2\) ，则 \(\left|A^{*} A A^{*}-6 A^{-1}\right|=\) \(\_\_\_\_\)．

\subsection{题目2}
2．若 1 是多项式 \(f(x)=2 x^{5}+a x^{2}-a x-2\) 的重根，则 \(a=\) \(\_\_\_\_\)．

\subsection{题目3}
3．设 \(A\) 的特征多项式为 \(\lambda^{3}-\lambda^{2}-\lambda+1\) ，则 \(A^{-1}=\) \(\_\_\_\_\)．

\subsection{题目4}
4．设 \(A\) 为3阶矩阵，它的特征值为－1，2，3，则 \(|A+2 E|=\) \(\_\_\_\_\)．

\subsection{题目5}
5．若 \(B, C\) 分别为 \(m\) 阶方阵和 \(n\) 阶方阵，则 \(\left|\begin{array}{ll}A & B \\ C & O\end{array}\right|=\)

\subsection{题目6}
6．设复矩阵 \(A\) 的不变因子为 \(1,1,1, \lambda-2,(\lambda-2)^{2},(\lambda-1)(\lambda-2)^{2}\) ，那么 \(A\) 属于特征值2的特征子空间的维数是 \(\_\_\_\_\)。

二．计算证明题．每题15分，共120分．\\

\subsection{题目7}
1．分别讨论 \(a, b\) 为何值时，线性方程组

\[
\left\{\begin{array}{l}
a x_{1}+x_{2}+x_{3}=1 \\
x_{1}+a x_{2}+x_{3}=-2 \\
x_{1}+x_{2}+a x_{3}=b
\end{array}\right.
\]

有解和无解？并在有无穷多解时求解。

\subsection{题目8}
2．设实二次型

\[
f\left(x_{1}, x_{2}, x_{3}\right)=2 x_{1}^{2}+5 x_{2}^{2}+2 x_{3}^{2}+4 x_{1} x_{2}-2 x_{1} x_{3}+2 a x_{2} x_{3}
\]

通过正交线性替换 \(X=Q Y\) 化为标准形 \(f=y_{1}^{2}+y_{2}^{2}+b y_{3}^{2}\) ．

\begin{itemize}
  \item[（1）] 求实数 \(a, b\) 。
  \item[（2）] 求正交矩阵 \(Q\) 。
\end{itemize}

\subsection{题目9}
3．设向量组

\[
\text { (I) : } \alpha_{1}=(1,2,-1)^{\mathrm{T}}, \alpha_{2}=(2,3,4)^{\mathrm{T}}, \alpha_{3}=(1,3, a)^{\mathrm{T}} ;
\]

\[
\text { (II) : } \beta_{1}=(1,3, b+4)^{\mathrm{T}}, \beta_{2}=(2,2, b+5)^{\mathrm{T}}, \beta_{3}=(3,1, b+7)^{\mathrm{T}} .
\]

求数 \(a, b\) 满足什么条件时向量组（I）与（II）等价。

\subsection{题目10}
4．设 \(V\) 是复数域上全体二阶矩阵构成的线性空间，定义 \(V\) 的一个线性变换 \(\sigma\) ：

\[
X \rightarrow\left(\begin{array}{ll}
1 & -2 \\
2 & -4
\end{array}\right) X, X \in V
\]

（1）求 \(\sigma\) 在基 \(E_{11}, E_{12}, E_{21}, E_{22}\) 下的矩阵，其中 \(E_{i j}\) 表示 \((i, j)\) 元素为 1 其余元素为 0 的矩阵。\\
（2）求 \(\sigma\) 的值域 \(\sigma(V)\) ，给出它的维数及一组基。\\
（3）求 \(\sigma\) 的核 \(N\) ，给出 \(N\) 的维数及一组基．

\subsection{题目11}
5．求复矩阵

\[
A=\left(\begin{array}{ccc}
0 & 2 & 3 \\
1 & 1 & 3 \\
-1 & -2 & -4
\end{array}\right)
\]

的若尔当标准形 \(J\) ，并求可逆矩阵 \(P\) ，使得 \(P^{-1} A P=J\) 。

\subsection{题目12}
6．设 \(V\) 是数域 \(P\) 上全体 \(n\) 阶对称矩阵关于矩阵加法，数乘运算构成的线性空间，\(U\) 是 \(V\) 中迹为 0的矩阵构成的子空间，\(W\) 是 \(V\) 中全体数量矩阵构成的子空间，分别求 \(U, W\) 的一组基，并证明

\[
V=U \oplus W .
\]

\subsection{题目13}
7．设 \(A\) 是 \(n\) 阶实对称矩阵，且 \(|A|=0, A\) 的前 \(n-1\) 个顺序主子式都是正数，证明：矩阵 \(A\) 是半正定的。

\subsection{题目14}
8．设 \(V\) 为数域 \(P\) 上的 \(n\) 维线性空间，线性变换 \(\mathcal{T}\) 在 \(V\) 的基 \(\varepsilon_{1}, \varepsilon_{2}, \cdots, \varepsilon_{n}\) 下的矩阵为

\[
\left(\begin{array}{llll}
a & & & \\
1 & a & & \\
& & & \\
& & 1 & a
\end{array}\right)
\]

其中 \(a \in P\) 。证明： \(\mathcal{T}\) 恰有 \(n+1\) 个不同的不变子空间。
\end{document}
